\section{Introduction}
\label{sec: Introduction}

Aerodynamic characterisation of Formula 1 cars is central to setup and race strategy. Teams measure $C_DA$ and $C_LA$ in wind tunnels and CFD but public access to these parameters is limited. The F1 live timing API, exposed through the FastF1 Python library \cite{FastF1}, provides car telemetry at 240 Hz and GPS position at 10 Hz for all sessions. This creates an opportunity to estimate aerodynamic parameters from first principles using freely available data.

Prior work on extracting aerodynamic parameters from road-car coast-down tests is well established, but F1 cars introduce several complications: active energy recovery during coasting (MGU-K), ride-height-dependent ground-effect aerodynamics, and the lack of direct brake force measurement. This report describes a methodology that addresses the energy recovery term explicitly. Sources of residual model mis-specification are investigated through gear-stratified and multi-driver sensitivity analyses; the approach is validated through cross-circuit comparison and internal consistency checks.

The Aston Martin AMR24 is used as the reference vehicle, selected by a data-driven survey of all 24 rounds and all 20 drivers in the 2024 season (\S\ref{sec: Constructor and Circuit Survey}). Five circuits spanning the full downforce range are analysed: Jeddah (low), Spa-Francorchamps (medium-low), Silverstone (medium), Yas Marina (medium-high), and Suzuka (high). All results use driver~14 (Alonso) throughout (see \S\ref{sec: Session Selection} for selection rationale).
